\chapter{Logica e calcolo: la risoluzione per la logica proposizionale}

\textit{Lucidi: \texttt{lucidi/4.Resolution\_PL}. Approfondimento:
Chang--Lee, cap.~4 (The Resolution Principle); Gabbrielli--Martini, capitolo
su LP.}

% =====================================================================
\section{Logica e calcolo}
% =====================================================================

\dfn{Logica}{
  Una \textbf{logica} è un linguaggio formale per esprimere enunciati, con:
  \begin{itemize}
    \item \textbf{sintassi}: le ``regole di ortografia'' delle espressioni;
    \item \textbf{semantica}: il significato delle espressioni (conseguenza
          logica $\models$);
    \item \textbf{calcolo}: un insieme di formule date e di regole
          \textbf{sintattiche} per manipolare le formule ed eseguire
          dimostrazioni. Una \textbf{derivazione} si scrive $\varphi \vdash
          \rho$ ($\varphi, \rho$ formule).
  \end{itemize}
}

\dfn{Calcolo}{
  Un calcolo è fatto di:
  \begin{itemize}
    \item \textbf{assiomi}: formule date, tautologie e contraddizioni
          elementari che non si possono derivare all'interno del calcolo;
    \item \textbf{regole di inferenza}: permettono di derivare nuove formule da
          formule date;
    \item \textbf{derivazione} $\varphi \vdash \psi$: una sequenza di
          applicazioni di regole di inferenza che parte dalla formula $\varphi$
          e termina con la formula $\psi$.
  \end{itemize}
}

$\models$ e $\vdash$ dovrebbero coincidere:
\begin{itemize}
  \item \textbf{Correttezza} (soundness): $\varphi \vdash \rho$ implica
        $\varphi \models \rho$ (ciò che si deriva è vero);
  \item \textbf{Completezza} (completeness): $\varphi \models \rho$ implica
        $\varphi \vdash \rho$ (ciò che è vero si deriva).
\end{itemize}

\dfn{Regola di inferenza}{
  Una regola di inferenza ha la forma
  \[ \frac{F_1, \dots, F_n}{F} \]
  \begin{itemize}
    \item le formule $F_1, \dots, F_n$ sono le \textbf{premesse};
    \item la formula $F$ è la \textbf{conclusione}.
  \end{itemize}
  Dato un insieme di formule $X$, se le premesse sono date (cioè contenute in
  $X$) allora la conclusione viene aggiunta a $X$.
  \begin{itemize}
    \item Un \textbf{passo di derivazione} $F_1, \dots, F_n \vdash F$ è
          l'applicazione di una regola di inferenza.
    \item Una \textbf{derivazione} è una sequenza di passi di derivazione, in
          cui le conclusioni vengono prese come premesse per i passi
          successivi.
  \end{itemize}
}

\begin{itemize}
  \item L'applicazione delle regole è di solito \textbf{non deterministica}
        (in ogni momento ci sono più regole applicabili, a più formule).
  \item L'insieme delle derivazioni si può rappresentare con \textbf{alberi di
        derivazione}.
\end{itemize}

\ex{La logica proposizionale vista finora}{
  \begin{itemize}
    \item \textbf{Sintassi}: simboli proposizionali (detti anche variabili
          proposizionali) e connettivi.
    \item \textbf{Semantica} $\models$: tavole di verità e assegnamento di
          valori di verità ai simboli proposizionali.
    \item \textbf{Calcolo} $\vdash$: la deduzione naturale.
      \begin{itemize}
        \item regole di inferenza: regole di introduzione ed eliminazione dei
              connettivi;
        \item derivazioni: alberi ottenuti usando regole e assiomi.
      \end{itemize}
  \end{itemize}
  Teorema di correttezza e completezza: $\varphi \vdash \rho$ sse
  $\varphi \models \rho$ (capitolo precedente).
}

\nt{
  La deduzione naturale è comoda per un umano ma scomoda per una macchina:
  ha molte regole, e a ogni passo bisogna scegliere quale applicare e
  (per $\wedge E$, $\to E$) \emph{indovinare} formule intermedie. La
  risoluzione nasce proprio per avere un calcolo facile da automatizzare.
}

% =====================================================================
\section{Il principio di risoluzione}
% =====================================================================

La risoluzione è stata introdotta da \textbf{Robinson (1965)}:
\begin{itemize}
  \item è una regola di inferenza \textbf{facile da implementare};
  \item è il meccanismo di esecuzione della programmazione logica (Prolog) e
        della programmazione logica \emph{con vincoli} (CLP);
  \item usa la \textbf{forma normale a clausole} (CNF) e, per la logica del
        primo ordine, l'\textbf{unificazione}.
\end{itemize}

\subsection{L'idea: refutazione}

Si lavora \textbf{per assurdo}, cioè per \textbf{refutazione}: la teoria
unita alla \textbf{negazione} della conseguenza deve essere
\textbf{insoddisfacibile}. Vogliamo dimostrare $\varphi \models \psi$; per il
teorema di refutazione (capitolo precedente)
\[
  \varphi \models \psi \quad\text{sse}\quad \varphi \wedge \neg\psi
  \ \text{è insoddisfacibile (inconsistente).}
\]

\nt{
  Sulla slide il teorema è scritto $\varphi \models \psi$ sse
  $\not\models \varphi \wedge \neg\psi$. Letto alla lettera,
  $\not\models F$ significa ``$F$ non è valida'', che è più debole (e qui
  sbagliato): $\varphi \wedge \neg\psi$ deve essere \textbf{falsa in tutte le
  interpretazioni}, non solo in almeno una. La versione giusta è quella del
  capitolo precedente: $\varphi \wedge \neg\psi$ \textbf{inconsistente},
  ovvero $\models \neg(\varphi \wedge \neg\psi)$.
}

\mprop{Il calcolo di risoluzione in tre passi}{
  Per dimostrare $\varphi \models \psi$:
  \begin{enumerate}
    \item[\textbf{1.}] trasformare tutte le formule (di $\varphi$ e
          $\neg\psi$) in formule equivalenti in \textbf{forma normale
          congiuntiva}, cioè in \textbf{clausole};
    \item[\textbf{2.}] applicare la \textbf{regola di risoluzione};
    \item[\textbf{3.}] fermarsi quando si deriva la \textbf{clausola vuota}
          (la contraddizione elementare, cioè \emph{false}).
  \end{enumerate}
}

\nt{
  Una formula in CNF $C_1 \wedge \dots \wedge C_k$ si tratta come
  l'\textbf{insieme di clausole} $\{C_1, \dots, C_k\}$ (la congiunzione è
  implicita), e ogni clausola come \textbf{insieme di letterali} (la
  disgiunzione è implicita). Così $A \vee A \vee B$ è la stessa clausola di
  $A \vee B$: i letterali ripetuti si fondono. Un insieme di formule
  $\{F_1, \dots, F_n\}$ diventa l'unione degli insiemi di clausole delle
  $F_i$.
}

\subsection{La regola di risoluzione}

\dfn{Calcolo di risoluzione per la logica proposizionale}{
  \textbf{Assioma}: la clausola vuota (cioè la contraddizione elementare).

  \textbf{Regola di risoluzione}:
  \[
    \frac{R \vee A \qquad R' \vee \neg A}{R \vee R'}
  \]
  dove $R$ e $R'$ sono disgiunzioni di letterali e $A$ è una formula
  atomica. La clausola $R \vee R'$ si chiama \textbf{risolvente}. La clausola
  vuota si indica di solito con $\square$.
}

\nt{
  \textbf{Come leggerla.}
  \begin{itemize}
    \item Si prendono due clausole che contengono una lo stesso atomo $A$
          positivo e l'altra negato (\textbf{letterali complementari}); si
          cancellano $A$ e $\neg A$ e si fa la disgiunzione di ciò che resta.
    \item Casi particolari già noti: con $R$ vuoto, cioè $\dfrac{A \quad
          \neg A \vee R'}{R'}$, è il \emph{modus ponens} ($A$ e $A \to R'$
          danno $R'$). Con $R$ e $R'$ entrambi vuoti, $\dfrac{A \quad \neg
          A}{\square}$: da una contraddizione si ottiene la clausola vuota.
    \item La clausola vuota è una disgiunzione di \emph{zero} letterali,
          quindi è \textbf{falsa in ogni interpretazione}: derivarla significa
          aver trovato una contraddizione.
    \item Sulla slide la clausola vuota è chiamata ``assioma''. È un uso un po'
          improprio: non è una formula di partenza ma il \emph{bersaglio}
          della ricerca, la contraddizione elementare a cui si vuole arrivare.
  \end{itemize}
}

\thm{Correttezza e completezza della risoluzione}{
  La risoluzione è \textbf{corretta e completa} per la logica proposizionale.
  Precisamente: un insieme di clausole $S$ è insoddisfacibile \textbf{sse}
  dalla chiusura di $S$ rispetto alla regola di risoluzione si deriva
  $\square$.
}

\pf{Correttezza: il risolvente è conseguenza logica delle premesse}{
  Sia $I$ un'interpretazione che rende vere $R \vee A$ e $R' \vee \neg A$.
  \begin{itemize}
    \item Se $A$ è vera in $I$, allora $\neg A$ è falsa, quindi deve essere
          vera $R'$; dunque $R \vee R'$ è vera.
    \item Se $A$ è falsa in $I$, allora deve essere vera $R$; dunque $R \vee
          R'$ è vera.
  \end{itemize}
  Quindi $R \vee A,\ R' \vee \neg A \models R \vee R'$. Per induzione sul
  numero di passi, ogni clausola derivata da $S$ è conseguenza logica di $S$.
  Se si deriva $\square$, che non ha modelli, allora $S$ non ha modelli: è
  insoddisfacibile.
}

\nt{
  \textbf{Attenzione: completa per la refutazione, non per la derivazione.}
  La risoluzione \emph{non} deriva tutte le conseguenze logiche: ad esempio
  $A \models A \vee B$, ma dalla sola clausola $A$ non si può applicare
  nessuna risoluzione, quindi $A \vee B$ non si deriva. Quello che vale è la
  \textbf{completezza per refutazione}: se l'insieme è insoddisfacibile,
  $\square$ si deriva. È per questo che si lavora sempre per assurdo:
  per mostrare $A \models A \vee B$ si considera $\{A, \neg A, \neg B\}$
  (la negazione di $A \vee B$ in clausole è $\neg A \wedge \neg B$), e da $A$
  e $\neg A$ si ottiene $\square$.
}

\nt{
  \textbf{Terminazione e decidibilità.} Con $n$ atomi esistono solo un numero
  finito di clausole distinte (ogni atomo compare positivo, negato o non
  compare: al più $3^n$ clausole). Quindi generando tutti i risolventi prima
  o poi non si ottiene più niente di nuovo (\textbf{saturazione}). A quel
  punto: se c'è $\square$ l'insieme è insoddisfacibile, altrimenti è
  soddisfacibile. La risoluzione proposizionale è quindi una
  \textbf{procedura di decisione}, coerente con la decidibilità vista nel
  capitolo precedente. (Nella logica del primo ordine invece la ricerca può
  non terminare.)
}

\subsection{Un esempio}

\ex{$\{A \to B,\ A\} \models B$}{
  \begin{enumerate}
    \item Trasformiamo le formule di $F = \{A \to B, A\}$ in forma a
          clausole: $F' = \{\neg A \vee B,\ A\}$.
    \item Aggiungiamo la negazione della conclusione che vogliamo dimostrare:
          $F'' = \{\neg A \vee B,\ A,\ \neg B\}$.
    \item Applichiamo la risoluzione:
          \begin{prooftree}
            \AxiomC{$\neg A \vee B$}
            \AxiomC{$A$}
            \BinaryInfC{$B$}
            \AxiomC{$\neg B$}
            \BinaryInfC{$\square$}
          \end{prooftree}
    \item Abbiamo ottenuto la clausola vuota: $F''$ è inconsistente, quindi
          $F \models B$.
  \end{enumerate}
}

\mprop{Perché la risoluzione e non la deduzione naturale}{
  Questo procedimento è molto più facile da automatizzare e implementare
  della deduzione naturale:
  \begin{itemize}
    \item \textbf{meno regole} (una sola);
    \item \textbf{meno non determinismo} (si sceglie solo quali due clausole
          risolvere, e su quale atomo; non bisogna inventare formule);
    \item inoltre (lo vedremo) la ricerca può essere \textbf{guidata dalla
          formula che vogliamo dimostrare}: è quello che fa Prolog, partendo
          sempre dal goal negato.
  \end{itemize}
}

% =====================================================================
\section{Esercizi svolti (approfondimento)}
% =====================================================================

\nt{
  Le slide contengono un solo esempio. Qui rifaccio con la risoluzione alcuni
  esercizi del capitolo precedente, risolti là con tavole e deduzione
  naturale. Lo schema è sempre lo stesso: (1) premesse in clausole,
  (2) negazione della tesi in clausole, (3) risolvere fino a $\square$.
}

\ex{Modus tollens (Chang--Lee 2.11)}{
  $P \to Q,\ \neg Q \models \neg P$?
  \begin{itemize}
    \item Premesse: $\neg P \vee Q$, \ $\neg Q$.
    \item Negazione della tesi: $\neg\neg P \equiv P$.
  \end{itemize}
  \begin{prooftree}
    \AxiomC{$\neg P \vee Q$}
    \AxiomC{$P$}
    \BinaryInfC{$Q$}
    \AxiomC{$\neg Q$}
    \BinaryInfC{$\square$}
  \end{prooftree}
  Quindi la conseguenza vale.
}

\ex{Esame 12/01/2023, es.\ 2 (prima parte)}{
  $(A \to B) \wedge (A \to C) \models A \wedge B \to C$?
  \begin{itemize}
    \item Premesse: $\neg A \vee B$, \ $\neg A \vee C$.
    \item Negazione della tesi: $\neg(A \wedge B \to C) \equiv \neg(\neg A
          \vee \neg B \vee C) \equiv A \wedge B \wedge \neg C$, cioè le tre
          clausole unitarie $A$, $B$, $\neg C$.
  \end{itemize}
  \begin{prooftree}
    \AxiomC{$\neg A \vee C$}
    \AxiomC{$A$}
    \BinaryInfC{$C$}
    \AxiomC{$\neg C$}
    \BinaryInfC{$\square$}
  \end{prooftree}
  La conseguenza vale. Notare che le clausole $\neg A \vee B$ e $B$ non
  servono: la risoluzione usa solo le clausole rilevanti.
}

\ex{Un insieme soddisfacibile: saturazione senza $\square$}{
  (Esame 12/01/2023, es.\ 3) $\{M \to C,\ \neg M \to T,\ C \vee T\}$ è
  contraddittorio? Clausole:
  \[
    C_1 = \neg M \vee C, \qquad C_2 = M \vee T, \qquad C_3 = C \vee T.
  \]
  L'unica coppia con letterali complementari è $C_1, C_2$ (su $M$), che dà
  $C \vee T$, cioè $C_3$: niente di nuovo. Le altre coppie non hanno
  letterali complementari. Siamo a \textbf{saturazione} senza aver ottenuto
  $\square$: l'insieme è \textbf{soddisfacibile} (come già trovato con la
  tavola, ad esempio $M = 0, C = 0, T = 1$).
}

\ex{Le lettere di Lewis Carroll}{
  Le nove premesse del capitolo precedente, in clausole:
  \begin{multicols}{3}
  \begin{enumerate}
    \item $\neg p \vee q$
    \item $s \vee \neg r$
    \item $\neg u \vee \neg t$
    \item $\neg v \vee p$
    \item $w \vee r$
    \item $\neg x \vee y$
    \item $\neg q \vee t$
    \item $v \vee \neg w$
    \item $\neg y \vee \neg s$
  \end{enumerate}
  \end{multicols}
  Tesi $x \to \neg u$; la sua negazione $\neg(\neg x \vee \neg u) \equiv x
  \wedge u$ dà le clausole \textbf{10.}~$x$ e \textbf{11.}~$u$.

  Partendo dalla tesi negata e risolvendo ogni volta con una premessa
  (\textbf{risoluzione lineare}):
  \begin{center}
  \begin{tabular}{rlcl}
    12. & $y$         & da & 10, 6 \\
    13. & $\neg s$    & da & 12, 9 \\
    14. & $\neg r$    & da & 13, 2 \\
    15. & $w$         & da & 14, 5 \\
    16. & $v$         & da & 15, 8 \\
    17. & $p$         & da & 16, 4 \\
    18. & $q$         & da & 17, 1 \\
    19. & $t$         & da & 18, 7 \\
    20. & $\neg u$    & da & 19, 3 \\
    21. & $\square$   & da & 20, 11
  \end{tabular}
  \end{center}
  L'argomento è valido. È la stessa catena di implicazioni del capitolo
  precedente, ma ottenuta in modo meccanico, senza dover ``vedere'' le
  contrapposizioni: le fa da sola la regola, perché nelle clausole
  $\neg q \vee t$ e $t \vee \neg q$ sono la stessa cosa.
}

\nt{
  \textbf{Errore classico: risolvere su due atomi alla volta.}
  Da $A \vee B$ e $\neg A \vee \neg B$ \textbf{non} si ottiene $\square$
  cancellando insieme $A/\neg A$ e $B/\neg B$: la regola elimina \emph{una
  sola} coppia complementare. I risolventi corretti sono $B \vee \neg B$ (su
  $A$) e $A \vee \neg A$ (su $B$), entrambi tautologie. Infatti l'insieme è
  soddisfacibile: $A = 1, B = 0$. Le clausole tautologiche (che contengono
  un letterale e il suo complementare) non servono mai a derivare $\square$ e
  si possono scartare.
}

\nt{
  \textbf{Legame con Davis--Putnam.} La \emph{unit propagation} dell'algoritmo
  di Davis--Putnam (capitolo precedente) è una risoluzione con una clausola
  unitaria: da $L$ e $\neg L \vee R$ si ottiene $R$, cioè ``si cancella $\neg
  L$ dalla clausola''. La ``null clause'' di Davis--Putnam è esattamente
  $\square$.
}
